\chapter{Einleitung}\label{ch:einleitung}


\section{Ausganssituation und Motivation}\label{sec:ausganssituation-und-motivation}

Zu nachgeführten Photovoltaikanlagen existieren bereits Arbeiten im Umfeld der HAW Hamburg und der TU Wien.
Als bisher bekannte fachliche Anknüpfungspunkte werden eine Bachelorarbeit Simulationen zur Ertragsoptimierung einer einachsig nachgeführten Photovoltaikanlage von Elena Steffens sowie eine Bachelorthesis zur Sonnennachführung von PV-Modulen von Franz Schlecht und eine Diplomarbeit zur Auslegung und Konstruktion einer einachsig nachgeführten Photovoltaik Anlage von Mathias Klotz genannt.
Diese Arbeiten zeigen, dass das Thema der Nachführung von Photovoltaikanlagen bereits wissenschaftlich und praktisch betrachtet wurde.
Für das vorliegende Projekt stehen jedoch nicht die allgemeine Auslegung neuer PV-Nachführsysteme, sondern die Restaurierung, Modernisierung und Erprobung eines bereits vorhandenen Modells im Mittelpunkt.

\section{Aufgabenstellung}\label{sec:aufgabenstellung}

Im Rahmen des Projekts ist die vorhandene PV-Nachführanlage technisch zu untersuchen, zu dokumentieren und soweit möglich wieder in einen funktionsfähigen Zustand zu versetzen.
Dazu sind zunächst der mechanische, elektrische und elektronische Aufbau sowie die bestehenden Schnittstellen durch Reverse Engineering zu erfassen.
Festgestellte Fehler und Schäden sind systematisch einzugrenzen und durch geeignete Reparatur- oder Modernisierungsmaßnahmen zu beheben.

Zusätzlich ist eine mikrocontrollerbasierte Steuerung zu entwickeln, die das PV-Modul automatisch nach der höchsten Lichteinstrahlung ausrichtet und eine manuelle Bedienung ermöglicht.
Die neue Steuerung soll zur vorhandenen elektrischen Schnittstelle kompatibel ausgeführt und separat montiert werden, sodass ein Wechsel zwischen der ursprünglichen und der neuen Steuerung möglich bleibt.
Abschließend sind die Funktionen der Gesamtanlage zu prüfen und die durchgeführten Arbeiten sowie die erzielten Ergebnisse nachvollziehbar zu dokumentieren.

\section{Zielsetzung}\label{sec:zielsetzung}

Ziel der Arbeit ist es, die vorhandene nachgeführte Photovoltaikanlage technisch zu untersuchen, zu dokumentieren und, soweit technisch möglich, zu restaurieren und zu modernisieren.

\begin{itemize}
    \item eine Projektplanung nach GPM-Vorlage,
    \item eine technische Bestandsaufnahme und Reverse-Engineering-Dokumentation
    \item eine soweit technisch mögliche Restaurierung des vorhandenen Modells,
    \item eine separat montierte, pin-kompatible Steuerung für den manuellen und automatischen Betrieb
    \item dokumentierte Funktions-, Bewegungs-, Genauigkeits- und Sicherheitstests,
    \item eine sinnvolle Kurzzeitmessung unter dokumentierten Bedingungen,
    \item eine Dokumentation für Lehr- und Demonstrationszwecke,
    \item eine abschließende Projektdokumentation einschließlich Schaltplänen und Programmcode.
\end{itemize}

\section{Abgrenzung}\label{sec:abgrenzung}

Die Arbeit umfasst sinnvolle Ertragsmessung über mehrere Tage.
Dies beinhaltet keine Simulationen oder Langzeitmessungen.
Die Tests werden bewusst kurzfristig geplant und umgesetzt, da sie abhängig von Wetter, Verfügbarkeit der Anlage und geeigneten Messbedingungen sind.
Außerdem ist nicht vorgesehen, eine vollständig neue kommerzielle PV-Nachführungsanlage zu entwickeln.
Der Schwerpunkt liegt auf der vorhandenen Anlage, ihrer technischen Wiederherstellung und ihrer Nutzbarkeit für Demonstrationszwecke, sowie der ausführlichen Dokumentation, die zusätzlich zu der Projektplanung stattfindet.

\section{Aufbau der Dokumentation}\label{sec:aufbau-der-dokumentation}

Die Dokumentation gliedert sich in einen projektorganisatorischen und einen technischen Teil.
Zunächst werden die Ausgangssituation, die Zielsetzung sowie die Planung und Durchführung des Projekts beschrieben.
Dazu gehören unter anderem die zeitliche Gliederung, die Aufgabenverteilung und der Umgang mit auftretenden Risiken.

Anschließend wird die technische Bearbeitung der PV-Nachführanlage dargestellt.
Diese umfasst die Untersuchung des Ist-Zustands, das Reverse Engineering, die Fehlersuche sowie die durchgeführten Reparatur- und Modernisierungsmaßnahmen.
Darauf aufbauend werden die Entwicklung, der Aufbau und die Funktionsweise der neuen Steuerung erläutert.
Abschließend folgen die Prüfung der Gesamtanlage, die Bewertung der Projektergebnisse und ein Ausblick auf mögliche weiterführende Arbeiten.

\subsection*{Transparenzhinweis zum Einsatz generativer KI}

Im Rahmen dieses Projekts wurde ChatGPT von OpenAI als unterstützendes Werkzeug eingesetzt.
Die Unterstützung umfasste insbesondere die Formulierung und sprachlichen Überarbeitung einzelner Dokumentationsabschnitte sowie zur Erstellung von \LaTeX{}-Beispielen.
Der Umgang mit generativer KI wurde bereits in der Lehrveranstaltung \enquote{Programmieren} im ersten Semester als mögliches Hilfsmittel für die Programmierung und die Bearbeitung derer Aufgaben thematisiert.
Der Einsatz in diesem Projekt erfolgte auf dieser Grundlage und wird an dieser Stelle offengelegt, um den Umfang der KI-Unterstützung transparent von den eigenen Entscheidungen, Bewertungen und praktischen Tätigkeiten abzugrenzen.

Der Einsatz generativer KI erfolgte, weil einige kleinere Nebenaufgabenbereiche ein hohes Maß an Zeit benötigen würden, und nicht viel zum Projekt beitragen würden.
Die KI-generierten Elemente werden an den jeweiligen Stellen gekennzeichnet und klar abgeschränkt.

Die Verantwortung für die Auswahl, Prüfung und Verwendung der übernommenen Inhalte sowie für die daraus gezogenen Schlussfolgerungen liegt bei der Projektgruppe.
Nicht ungeprüft KI-generierte Ausgaben werden nicht als fachlich korrekt vorausgesetzt oder als Fakten angegeben.


